\documentclass[conference]{IEEEtran}
\usepackage[T1]{fontenc}
\usepackage[utf8]{inputenc}
\usepackage{amsmath,amssymb}
\usepackage{array}
\usepackage{booktabs}
\usepackage{multirow}
\usepackage{url}
\usepackage{microtype}
\newcolumntype{P}[1]{>{\raggedright\arraybackslash}p{#1}}
\setlength{\emergencystretch}{2em}

\title{What Creates the Earth's Magnetic Field?\\A Coupled Core-to-Space Model for Field Motion}
\author{Research Plan Author Agent\\Xcientist four-agent simulation}

\begin{document}
\maketitle

\begin{abstract}
Earth's magnetic field is not produced by rotation alone and is not a rigid dipole that simply drifts through space. The leading physical picture is a self-sustaining geodynamo: electrically conducting liquid metal in the outer core moves under thermal and compositional buoyancy, rotation organizes the flow, induction amplifies magnetic structure, and magnetic diffusion and Lorentz feedback close the system. The field observed at Earth's surface is additionally altered by lithospheric, oceanic, ionospheric, and magnetospheric currents. This report proposes GEODYNAMO-X, a coupled core-to-space state-space and data-assimilation experiment designed to explain why the field changes and to separate internal secular variation from external disturbances. The Survey stage identifies core-flow non-uniqueness, source mixing, multi-scale coupling, and weak forecast validation as the central gaps. The Idea stage selects a model with latent core magnetic and flow states, conductivity transfer operators, explicit external-current states, and physics-constrained residual learning. The ExperimentDesign stage specifies induction and rotating-flow equations, synthetic-twin identifiability tests, observatory and satellite observation operators, rolling-origin forecasts, baselines, ablations, and falsification rules. This is a design-only proposal: it reports no completed reconstruction, pole trajectory, forecast score, or new geophysical measurement.
\end{abstract}

\begin{IEEEkeywords}
geodynamo, geomagnetic field, secular variation, core flow, magnetic pole motion, ionosphere, magnetosphere, data assimilation
\end{IEEEkeywords}

\section{Introduction}
The Earth's magnetic field is often summarized by saying that the planet's rotation generates it. That statement identifies an important organizing influence but omits the energy source, the conducting fluid, the induction process, and the distinction between internal and external signals. The field is best posed as a coupled dynamical and inverse problem: what flow in the electrically conducting outer core can sustain the long-lived internal field, and which parts of an observed field change originate in the core rather than in the ionosphere, magnetosphere, crust, ocean, or instrument?

The scientific reframe used in this report is:

\begin{quote}
How do thermal and compositional convection in Earth's conducting outer core, organized by rotation and constrained by the core--mantle boundary, sustain the main field; and how can internal dynamics be separated from external and induced sources to explain and predict observed field motion?
\end{quote}

This reframe distinguishes field generation from field observation. A conducting fluid moving through a magnetic field can induce current and magnetic field. Rotation constrains the geometry of the flow through the Coriolis force, while buoyancy supplies mechanical energy and magnetic induction competes with diffusion. At the surface, the internal field is superposed with currents in the ionosphere and magnetosphere and with induced and lithospheric contributions. The magnetic poles therefore move in the sense that the locations inferred from the internal field change; no rigid magnetic object is required to travel through the mantle.

The physical mechanism class is strongly supported by geodynamo theory and numerical dynamo work, but the Earth-specific inverse problem remains open. Reviews of Earth's magnetism emphasize the interaction of core convection, rotation, electrical induction, and dissipation~\cite{robertsking,buffett}. Three-dimensional simulations show that a self-sustaining and reversing field can arise in rotating spherical shells~\cite{glatzmaier}. Scaling studies connect numerical dynamos to planetary-field observables while also exposing the gap between accessible simulations and Earth's extreme parameter regime~\cite{christensen}. Community geomagnetic models such as the International Geomagnetic Reference Field (IGRF) organize global observations into main-field and secular-variation representations, but a reference field is not by itself a unique causal model or a validated future forecast~\cite{alken}.

This report makes three contributions. First, it turns the broad question into a source-aware test of field generation and motion. Second, it proposes GEODYNAMO-X, which assimilates ground and satellite observations into a constrained latent-state model containing both core and external sources. Third, it defines a synthetic-twin gate, held-out forecast protocol, and falsification criteria so that improved fit cannot be confused with recovered physics.

\section{Survey: Physical Evidence and Open Gaps}
\subsection{The internal geodynamo}
The outer core is a liquid, electrically conducting region. Secular cooling and the growth of the solid inner core can generate thermal and compositional buoyancy. Fluid motion transports magnetic flux, inducing electric currents that generate a field. The system remains active only if induction and available mechanical power offset magnetic diffusion and Ohmic dissipation. This is why the phrase “energy made of rotation” is incomplete: rotation organizes the motion but does not replace buoyancy or induction.

The Coriolis force favors rapidly rotating, columnar flow structures and influences helicity and the conversion between kinetic and magnetic energy. The Lorentz force feeds back on the flow, so the magnetic field is not a passive tracer. Pressure, viscosity, compressibility, boundary heat flux, inner-core morphology, electrical conductivity, and unresolved turbulence influence the solution. Heterogeneity at the core--mantle boundary can imprint preferred spatial structures, but the inverse mapping from a surface field to the boundary flow is non-unique.

Numerical dynamos are therefore mechanism tests and prior generators. They can demonstrate self-sustaining induction, polarity reversals, and secular variation in a physically constrained system. They do not establish that one simulation has the same heat flux, diffusivity, turbulence, or inner-core state as Earth. A responsible model must carry structural uncertainty instead of selecting a single simulation as ground truth.

\subsection{Internal, external, and induced field motion}
On years-to-centuries time scales, the main internal field varies because the core magnetic field and its source flow evolve. The time derivative of the internal spherical-harmonic coefficients is called secular variation. It changes field intensity and direction and therefore changes the location of a modeled magnetic pole. Such motion is a projection of a changing field, not evidence that the magnetic axis is a solid rod.

On seconds-to-days time scales, ionospheric and magnetospheric currents respond to solar-wind forcing, auroral electrodynamics, tides, conductivity, and local time. These currents can generate large transient magnetic signals. They can also induce currents in the conducting Earth and in technological infrastructure. A surface observation during a geomagnetic storm is not a clean measurement of core evolution. Source separation must therefore be part of the model or its error must be explicitly quantified.

The observation problem is multi-scale. Ground observatories offer long records at fixed locations. Satellites offer global coverage and changing geometry, with simultaneous position, attitude, and space-environment information. Paleomagnetic records extend the time window to reversals and excursions but are sparse and filtered by geological recording processes. IGRF and related models provide a common harmonic representation, while Swarm-era studies demonstrate the value of combining satellites with ground observatories for time-varying field estimation~\cite{alken,finlay}.

\begin{table*}[t]
\centering
\caption{Evidence boundaries used by the proposed GEODYNAMO-X study.}
\label{tab:evidence}
\small
\begin{tabular}{P{2.7cm}P{4.3cm}P{5.0cm}P{3.0cm}}
\toprule
Evidence layer & Supported physical proposition & Use in the design & Inference not permitted \\
\midrule
Geodynamo theory & Conducting-fluid motion, buoyancy, rotation, induction, diffusion, and Lorentz feedback form a coupled mechanism & Define constrained transition priors and energy-accounting diagnostics & Rotation alone is the field source \\
Numerical dynamos & Rotating spherical-shell systems can sustain and reverse magnetic fields in a mechanism-compatible regime & Generate synthetic twins and structural prior families & One simulation is a unique Earth solution \\
Geomagnetic models & Global observations can be represented by internal-field coefficients and secular variation & Define common observables and forecast targets & A fitted reference field is a causal forecast \\
Observatory/satellite records & Field variation contains information about internal and external current systems & Build source-aware observation operators and holdouts & Every rapid residual is core flow \\
Paleomagnetism & Long records constrain long-term direction, intensity, excursions, and reversals & Test time-scale consistency with a separate recording model & Sparse records identify instantaneous flow \\
\bottomrule
\end{tabular}
\end{table*}

\subsection{Survey gap ledger}
The first gap is energy-to-field closure. Available heat and composition budgets are uncertain, while numerical models cannot resolve the full diffusivity and turbulence hierarchy. The second is core-flow non-uniqueness: surface secular variation underdetermines three-dimensional flow at the core--mantle boundary. The third is internal--external source mixing, especially during disturbed intervals. The fourth is multi-scale coupling; core, mantle, ocean, ionosphere, and magnetosphere are often represented in separate stages even though their signals meet at the observation operator.

The fifth gap is time-scale transfer. Instrumental records are short compared with reversal histories, while paleomagnetic records are long but irregular. The sixth is forecast validation. Models can fit historical coefficients without forecasting held-out vectors or uncertainty. The seventh is parameter identifiability. Rotation, buoyancy, conductivity, boundary heat flow, and unresolved closure terms can compensate for one another in an inverse model. These gaps motivate a design in which recovery of known latent states precedes interpretation of real records.

\subsection{Time scales and observation geometry}
The field-motion question is best analyzed as a hierarchy rather than a single time series. At the shortest relevant scales, external currents driven by solar-wind and auroral activity create rapid disturbances. At intermediate scales, ionospheric tides, seasonal conductivity, and recurrent forcing modulate the external field. On annual to centennial scales, the core generates secular variation that changes the internal field coefficients. On millennial and longer scales, excursions and reversals constrain the history of the dynamo. The same sensor can record several of these processes at once, so the time scale of an observation does not identify its source.

Spatial geometry creates a second hierarchy. High-degree lithospheric structure is strongest near the surface, while the large-scale core field is smoother because it passes through the mantle. External current systems have characteristic local-time, latitude, and solar-wind dependencies. Satellite altitude and track direction determine which mixtures are visible; fixed observatories provide temporal continuity but cannot alone resolve global spatial modes. GEODYNAMO-X therefore uses sensor geometry in the observation operator rather than treating all field vectors as exchangeable rows.

The field representation is a spherical-harmonic expansion whose coefficients are indexed by degree, order, and epoch. Low-degree internal coefficients are the main target for core secular variation, while higher-degree terms and regional transfer effects are retained as nuisance or separate latent components. The harmonic truncation is chosen before evaluation and varied in sensitivity analysis. A change in a low-degree coefficient can move a derived pole substantially, whereas a local high-degree anomaly may have little effect on the pole. Reporting both coefficient uncertainty and pole uncertainty prevents a visually dramatic local anomaly from being misread as a global axis motion.

\begin{table}[t]
\centering
\caption{Source classes, characteristic observables, and validation emphasis.}
\label{tab:timescales}
\footnotesize
\begin{tabular}{P{1.45cm}P{1.65cm}P{2.8cm}}
\toprule
Source class & Typical time scale & Validation emphasis \\
\midrule
Core & Years to millennia & Secular variation, internal coefficients, flow recovery \\
Ionosphere & Seconds to days; seasonal & Local time, latitude, forcing response \\
Magneto\\sphere & Seconds to days & Storm and substorm event holdouts \\
Induced layers & Minutes to years & Conductivity transfer and regional phase \\
Lithosphere & Static to geological & Spatial residual and harmonic separation \\
Paleo-\\magnetic & Centuries to millions of years & Separate recording and smoothing model \\
\bottomrule
\end{tabular}
\end{table}

The table does not assign a hard cutoff to any source. Its purpose is to define what evidence can challenge each layer. For example, an ionospheric model should be tested during changes in local time and solar forcing, while a core-flow model should be tested on secular-variation intervals and spatially withheld sensors. A model that achieves a low all-period error by smoothing away storms or by fitting lithospheric structure is not a successful explanation of core dynamics.

The observation geometry also controls identifiability. If all sensors sample a narrow region or a single temporal regime, an internal and external source may be observationally equivalent. Adding a satellite track with a different altitude or local-time distribution can break that equivalence, but only if its calibration and forcing metadata are retained. The synthetic twin explicitly varies these geometries and reports the posterior contraction obtained from each observation class. This gives a measurable answer to the question of what additional measurement would most reduce uncertainty, instead of assuming that more data are automatically informative.

The design also distinguishes interpolation from extrapolation. Interpolation tests whether a model can reconcile multiple sensors and source layers within an observed interval. Extrapolation tests whether the latent dynamics produce a useful future distribution after the last assimilated observation. The former is necessary for a coherent field map; the latter is required for a forecast claim. A model can pass the first while failing the second, particularly when its learned residual captures historical instrument artifacts. All report tables therefore carry an explicit label for reconstruction, interpolation, or forecast.

The survey rejects five tempting overclaims. Rotation does not create a magnetic field without conducting-fluid motion and energy. Ionospheric currents are not the primary source of the long-lived internal dipole. A pole track is not a literal trajectory of a rigid object. A numerical dynamo is not an observation of Earth's core. Finally, a good historical fit is not a validated forecast. These boundaries permit a strong physical hypothesis while keeping the claims testable.

\section{Idea: GEODYNAMO-X}
\subsection{Competing routes}
The core-only inversion route is interpretable and directly targets secular variation, but it can absorb external-current residuals into an artificial core acceleration. Static-reference and linear secular-variation extrapolations are transparent engineering baselines, but they do not represent a changing source mechanism. An uncoupled internal-plus-external regression tests source separation but cannot express shared boundary conditions or feedback. A numerical-dynamo ensemble supplies a physical prior but may not be identifiable or observationally calibrated.

GEODYNAMO-X combines the useful parts of these routes. It represents the core magnetic field and low-dimensional boundary flow as latent states; uses conductivity operators for mantle, lithosphere, and ocean transfer; includes explicit ionospheric and magnetospheric states; and maps all sources to the measurement geometry of each observatory or satellite. A constrained residual model may correct unresolved closure error, but it is prohibited from creating magnetic monopoles, violating the induction constraint, or relabeling an external pulse as core flow without an attribution penalty.

\subsection{Hypothesis and predictions}
The central hypothesis is:

\begin{quote}
A coupled core-to-space state-space model with physical priors and source-specific observation operators will produce better calibrated held-out predictions of vector-field change and internal-field pole motion than static extrapolation, core-only inversion, or uncoupled source separation, while recovering known latent states in synthetic twins.
\end{quote}

The first prediction is a reduction in held-out vector-field and secular-variation error when joint assimilation is used. The second is a reduction in false core-flow acceleration during intervals dominated by external current changes. The third is that physical priors improve extrapolation more than in-sample fit. The fourth is that synthetic twins will reveal which flow scales and buoyancy parameters are identifiable; unresolved scales should remain uncertain rather than receive spurious precision.

The direction is falsifiable. If the full model does not improve held-out skill under equal data and model-budget accounting, the coupling is not justified for that domain. If it predicts well but fails source recovery, it remains a reconstruction tool but does not support the causal claim. If a simpler model matches it under all declared tests, the simpler model is preferred.

\section{ExperimentDesign}
\subsection{State-space formulation}
The proposed system has four layers: an outer-core magnetic and flow state, a conducting transfer layer, an external-current state, and an instrument state. Let the latent state contain harmonic coefficients and reduced flow modes. The core magnetic evolution is represented by:

\begin{equation}
\frac{\partial \mathbf{B}}{\partial t}=\nabla\times(\mathbf{u}\times\mathbf{B})+\eta\nabla^2\mathbf{B},\qquad \nabla\cdot\mathbf{B}=0,
\label{eq:induction}
\end{equation}

where $\mathbf{B}$ is magnetic flux density, $\mathbf{u}$ is outer-core fluid velocity, and $\eta$ is magnetic diffusivity. The divergence-free condition is enforced in the basis representation. The first term is field induction by flow and the second is magnetic diffusion.

The rotating-flow prior is:

\begin{equation}
\rho\left(\frac{\partial \mathbf{u}}{\partial t}+\mathbf{u}\cdot\nabla\mathbf{u}+2\boldsymbol{\Omega}\times\mathbf{u}\right)=-\nabla p+\rho'\mathbf{g}+\mathbf{J}\times\mathbf{B}+\nabla\cdot\boldsymbol{\tau},
\label{eq:momentum}
\end{equation}

where $\rho$ is reference density, $\boldsymbol{\Omega}$ is planetary rotation, $p$ is pressure, $\rho'\mathbf{g}$ is buoyancy, $\mathbf{J}$ is electric current density, and $\boldsymbol{\tau}$ contains viscous and unresolved stresses. Rotation enters as an organizing force; it is not an independent magnetic source.

The discrete transition is:

\begin{equation}
\mathbf{x}_{t+1}=\mathcal{F}_{\theta}(\mathbf{x}_t,\mathbf{q}_t,\mathbf{w}_t)+\boldsymbol{\epsilon}_t,
\label{eq:transition}
\end{equation}

where $\mathbf{x}_t$ is the full latent state, $\mathbf{q}_t$ contains forcing and boundary covariates, $\mathcal{F}_{\theta}$ is the constrained transition, $\mathbf{w}_t$ is a learned closure residual, and $\boldsymbol{\epsilon}_t$ is process uncertainty. The residual is regularized by divergence, energy, and source-attribution constraints.

The measurement equation is:

\begin{equation}
\mathbf{y}_{k,t}=\mathcal{H}_{k,t}(\mathbf{x}_t)+\mathbf{b}_{k,t}+\boldsymbol{\nu}_{k,t},
\label{eq:observation}
\end{equation}

where $\mathbf{y}_{k,t}$ is the vector measurement from sensor $k$, $\mathcal{H}_{k,t}$ maps all source layers to the sensor's position and attitude, $\mathbf{b}_{k,t}$ is calibration bias, and $\boldsymbol{\nu}_{k,t}$ is measurement noise. Different ground and satellite geometries share a latent source state but retain separate observation operators.

The surface signal is reported as:

\begin{equation}
\mathbf{B}_{\mathrm{obs}}=\mathbf{B}_{\mathrm{core}}+\mathbf{B}_{\mathrm{ext}}+\mathbf{B}_{\mathrm{ind}}+\mathbf{B}_{\mathrm{lith}}+\boldsymbol{\delta}_{\mathrm{inst}},
\label{eq:source}
\end{equation}

where the terms denote core, external, induced, lithospheric, and instrument contributions. The equation is a source-accounting interface, not a claim that all terms are equally identifiable at every time and location.

\subsection{Data inventory and splitting}
The proposed inventory contains vector measurements from geomagnetic observatories, satellite magnetic measurements with position and attitude data, solar-wind and geomagnetic-activity covariates, global reference-field and secular-variation coefficients, and selected paleomagnetic records. Every record receives a coordinate frame, time standard, calibration flag, uncertainty, source label, and availability interval.

Preprocessing is frozen before training. It includes coordinate conversion, quality flags, common time bins, known spike removal, treatment of missingness, and a predeclared rule for quiet and disturbed intervals. Future external covariates are not made available to a historical forecast. Paleomagnetic observations use a separate geological recording operator and are not treated as dense instrumental samples.

The temporal split uses rolling origins with a training interval, assimilation interval, and forecast horizon. Spatial holdouts reserve observatories or satellite tracks. Event holdouts reserve selected disturbance intervals. The final evaluation mask is locked before hyperparameter selection. This prevents a model from learning the forecast target through accidental overlap in time, sensor, or external forcing.

\begin{table}[t]
\centering
\caption{Model comparisons and their scientific roles.}
\label{tab:baselines}
\footnotesize
\begin{tabular}{P{1.95cm}P{1.95cm}P{2.15cm}}
\toprule
Model & Role & Main limitation tested \\
\midrule
Static reference & Engineering baseline & No source dynamics \\
Linear variation & Short-horizon baseline & Trend extrapolation only \\
Core-only & Internal inversion baseline & External residual leakage \\
Uncoupled split & Source-separation baseline & No shared dynamics \\
Dynamo prior & Physical-prior baseline & No observational assimilation \\
GEODYNAMO-X & Primary treatment & Coupling and identifiability cost \\
Oracle source & Synthetic upper bound & Not deployable \\
\bottomrule
\end{tabular}
\end{table}

\subsection{Synthetic-twin identifiability gate}
Before historical fitting, a numerical-dynamo ensemble generates hidden core fields and flow states. A conductivity-transfer model adds mantle, lithosphere, and ocean effects. An external-current generator creates quiet, storm, and missing-forcing regimes. Ground and satellite observation operators sample the result with realistic geometry, noise, gaps, and calibration shifts. The inference model receives only the observations and covariates, not the latent truth.

The gate measures recovery of internal-field coefficients, core-flow modes, external-current amplitudes, pole location, and predictive uncertainty. It varies sensor density, satellite-track coverage, forcing availability, harmonic truncation, and structural model discrepancy. A parameter is considered identifiable only if posterior error shrinks with informative data, coverage remains calibrated, and recovery is stable across plausible model families.

The synthetic twin includes an intentionally difficult case: a rapid external disturbance resembles a regional core secular-variation pulse. GEODYNAMO-X must assign the pulse to external or uncertain sources rather than automatically accelerating the core. A small total-field error does not excuse a wrong causal attribution.

\subsection{Historical reconstruction and forecast}
After passing the synthetic gate, all models are fit to historical-like records under the locked splits. The retrospective task estimates the field and its decomposition. The forecast task predicts future vector field and secular variation at short, intermediate, and decadal horizons. The internal-field pole location is derived from the inferred internal coefficients and reported with an uncertainty ellipse; it is not directly optimized as a label.

Primary metrics are vector-field root-mean-square error, angular error, secular-variation error, pole-location error, negative log predictive density, interval coverage, and synthetic source-attribution error. Secondary metrics measure degradation under missing sensors, calibration drift, storm contamination, spatial holdouts, and changed satellite coverage. Skill is reported relative to static and linear baselines across time blocks, not as one average that hides failure during important events.

\subsection{Ablations and sensitivity}
Eight ablations isolate the mechanism: remove the Coriolis prior; remove compositional-buoyancy covariates; replace heterogeneous core--mantle boundary heat flow by a uniform boundary; remove external-current states; replace the conductivity transfer by an identity operator; remove divergence-free and energy-accounting residual constraints; reduce observatory or satellite coverage; and remove uncertainty propagation. Each ablation is evaluated with the same data split and forecast protocol.

Sensitivity analysis varies magnetic diffusivity, flow-mode rank, harmonic degree, external-forcing lag, observation noise, source correlation, forecast horizon, and model-discrepancy family. Nested sampling separates structural uncertainty from parameter uncertainty. If a qualitative conclusion changes under a plausible structure, the output is labeled model-dependent. No ensemble average is permitted to erase a sign change.

\section{Expected Interpretations and Falsification}
Four outcome classes are preregistered. A \textit{coupled-source gain} occurs if GEODYNAMO-X improves held-out field and attribution skill with calibrated uncertainty. A \textit{physical-prior gain} occurs if the physics improves extrapolation but not reconstruction, indicating value mainly at longer horizons. A \textit{separation-only gain} occurs if external states help while the core dynamics add no held-out benefit. A \textit{no-coupling gain} occurs if a simpler model matches or exceeds the full model under equal constraints.

\begin{table}[t]
\centering
\caption{Predeclared interpretation and falsification rules.}
\label{tab:rules}
\footnotesize
\begin{tabular}{P{1.55cm}P{2.35cm}P{2.25cm}}
\toprule
Outcome & Required signature & Meaning \\
\midrule
Coupled gain & Held-out skill and attribution improve; intervals calibrated & Coupling supported \\
Prior gain & Extrapolation improves; reconstruction does not & Physics helps forecast \\
Separation gain & External states improve attribution only & Source mixing is binding \\
No-coupling & Simple baseline matches or wins & Full model not justified \\
\bottomrule
\end{tabular}
\end{table}

GEODYNAMO-X is falsified if it fails synthetic-twin state recovery, falsely assigns the adversarial external pulse to core acceleration, produces systematically overconfident intervals, or loses its held-out advantage when baselines receive identical data, horizon, and parameter-budget treatment. It is also falsified as a causal model if multiple incompatible source decompositions fit equally well and the posterior cannot express that ambiguity.

A negative result is scientifically useful. If external states alone improve forecasts, the field-motion problem is primarily a source-separation bottleneck for the tested horizon. If the physical prior improves only long-horizon extrapolation, the model informs future uncertainty without claiming detailed core-flow recovery. If the static model wins, added dynamical complexity is not justified for that operational setting. None of these outcomes establishes that a reversal will occur or that a magnetic pole will follow a predetermined path.

\section{Data Contract and Evaluation Protocol}
The experiment cannot establish a causal geodynamo claim unless its observations and inferred quantities remain distinguishable. Before fitting any model, each record is assigned a stable identifier, time standard, coordinate frame, sensor geometry, calibration state, quality flag, uncertainty estimate, and missingness interval. Observatory measurements are tagged with site metadata and local environmental flags. Satellite measurements retain position, attitude, sampling, and external-forcing context. Reference-field coefficients are stored with their harmonic degree, epoch, and secular-variation convention. Paleomagnetic records are tagged with their geological smoothing and recording uncertainty rather than inserted as if they were dense modern observations.

The data contract separates five objects. A \textit{measured field} is an instrument output after declared calibration. An \textit{inferred internal field} is a posterior quantity attributed to the core through the source model. An \textit{external field} is a posterior quantity associated with ionospheric or magnetospheric current modes. An \textit{induced field} is the response of conducting layers to time-varying sources. A \textit{forecast field} is a prediction made without access to the future target. These names are retained in every output file. A pole location is stored with the internal-field definition and harmonic truncation that generated it.

The experiment uses rolling temporal origins. For each origin, the training window estimates parameters, the assimilation window updates the latent state, and the forecast window is sealed until prediction is written. Spatial holdouts reserve observatory sites or satellite tracks. Event holdouts reserve intervals with strong external forcing. Hyperparameters are selected only within the development split. The final temporal, spatial, and event masks are never used to tune the model. This design directly tests whether the model extrapolates rather than reconstructs information that leaked across a boundary.

The primary score is a vector-field error paired with a calibration score. Let $\widehat{\mathbf{B}}_{k,t}$ be the predictive mean and $\mathbf{B}_{k,t}$ the held-out vector. The vector loss is:

\begin{equation}
L_{\mathrm{vec}}=\frac{1}{N}\sum_{k,t}\left\|\widehat{\mathbf{B}}_{k,t}-\mathbf{B}_{k,t}\right\|_2^2,
\label{eq:vectorloss}
\end{equation}

where $N$ is the number of held-out sensor-time pairs and $\|\cdot\|_2$ is Euclidean vector norm. This metric treats intensity and direction jointly. Angular error, secular-variation error, and derived pole-location error are reported separately so that a good total vector score cannot conceal a directional failure.

Predictive distributions are evaluated with negative log predictive density and empirical coverage of declared intervals. Coverage is measured separately for quiet intervals, disturbed intervals, spatial holdouts, and forecast horizons. A model that reduces RMSE by producing overconfident intervals is not accepted as an improvement. For synthetic twins, source-attribution error is the fraction of generated external or core pulses assigned to the wrong latent source. For inferred core-flow modes, posterior rank and recovery error are reported, not only the posterior mean.

The decision table is deliberately multi-dimensional. The primary treatment must be compared on field error, secular variation, pole uncertainty, source attribution, and calibration. Secondary reports include compute cost, sensitivity to missing sensors, degradation under calibration drift, and stability across structural transition families. The analysis records the first failure mode: wrong source attribution, non-identifiability, forecast drift, interval undercoverage, or observation-operator mismatch. This makes the experiment useful even when the headline hypothesis fails.

\begin{table}[t]
\centering
\caption{Minimum data and output contract for an interpretable run.}
\label{tab:contract}
\footnotesize
\begin{tabular}{P{1.7cm}P{2.15cm}P{2.2cm}}
\toprule
Object & Required metadata & Prohibited shortcut \\
\midrule
Measurement & Time, frame, geometry, calibration, uncertainty & Treat missing data as zero \\
Internal field & Source label, harmonic degree, posterior & Call a fit a direct observation \\
External field & Forcing context, current mode, uncertainty & Absorb storms into core flow \\
Pole position & Derived definition and confidence region & Optimize pole independently \\
Forecast & Origin, horizon, mask, covariates & Use future targets or forcing \\
\bottomrule
\end{tabular}
\end{table}

The final acceptance rule is conjunctive. GEODYNAMO-X must pass synthetic state recovery, improve held-out calibrated prediction over the strongest baseline, and avoid systematic source-attribution errors under the adversarial external pulse. A gain that appears only after removing hard observations, relaxing the parameter budget, or changing the target definition is not considered evidence for the full model. This protocol keeps the report ambitious about a unified core-to-space explanation while making every claimed improvement traceable to a declared data and evaluation choice.

\section{Implementation, Governance, and Reproducibility}
The computational stack contains a data registry, source-specific observation operators, constrained state transitions, a transfer-layer module, a probabilistic assimilation engine, forecast evaluation, and an audit layer. Each observation, coordinate transform, time split, parameter prior, model version, random seed, checkpoint, and evaluation mask receives an identifier. Changes to the source registry create a new run rather than silently rewriting a prior result.

The registry distinguishes measured field, inferred internal field, inferred external field, and forecast field. It stores uncertainty and missingness rather than imputing a perfect record. A forecast is not promoted to an observation. A pole position is labeled by the harmonic truncation and internal-field definition used to derive it. This naming discipline makes it possible to compare forecasts without confusing a model artifact with a physical measurement.

Operational geomagnetic services remain the authority for safety-relevant warnings. GEODYNAMO-X, if ever evaluated beyond this design, would be compared with established services and would expose uncertainty, source attribution, and failure modes. It would not issue an automatic warning solely because a latent core state crossed a model threshold. The present work ends before deployment and contains no operational alert.

The proposal includes no empirical result. The expected-results fields remain empty until the synthetic and historical experiments are actually run. Reproducibility is achieved through versioned configuration, locked data splits, synthetic truth, saved masks, and complete parameter provenance rather than by claiming that the equations alone determine a unique answer.

\section{Risks and Scope}
The largest scientific risk is model discrepancy. Numerical dynamos, core-flow inversions, conductivity transfer models, external-current models, and sensor calibration models may disagree. GEODYNAMO-X treats this disagreement as a structured uncertainty rather than allowing a flexible residual to hide it. A recommendation that survives only one transition family is not robust.

The largest inverse-problem risk is non-identifiability. Surface observations constrain only certain spatial and temporal modes of the core. The report therefore treats a posterior family of flows as a valid result. A broad uncertainty interval is preferable to a precise but non-identifiable flow map.

The largest attribution risk is external contamination. Solar-wind and ionospheric forcing can create rapid changes that resemble regional internal variation. Event holdouts and the adversarial synthetic twin target this failure directly. The largest communication risk is to describe pole motion as if it were a rigid object moving through Earth; the report instead defines pole location as a derived property of the internal field.

This report is a research proposal, not a new geomagnetic model, a forecast, a reversal prediction, or an empirical estimate of core-flow speed. Its direct scientific claim is narrower: a conducting-fluid geodynamo provides the physical source class, rotation organizes but does not alone generate the field, and observed field motion should be tested with a coupled model that separates internal secular variation from external electromagnetic sources.

\section{Conclusion}
The Earth's long-lived magnetic field is most plausibly sustained by induction in a moving, electrically conducting liquid outer core. Thermal and compositional buoyancy provide energy, rotation organizes the flow, magnetic diffusion removes structure, and Lorentz feedback couples the field back to the fluid. The field moves because the core state changes and because external current systems vary; the apparent motion of a magnetic pole is a derived change in the internal field, not a rigid axis sliding through the planet.

GEODYNAMO-X makes this explanation testable. It combines a constrained induction model, rotating-flow priors, conductivity transfer, explicit ionospheric and magnetospheric states, and sensor-specific observation operators. It first asks whether the latent sources can be recovered in synthetic twins, then tests held-out vector-field forecasts and pole-location uncertainty under real-data-like sampling. Baselines, ablations, structural sensitivity, and falsification rules prevent a complex model from winning merely by fitting history.

The answer to the motivating question is therefore direct but qualified: the field is created primarily by a self-sustaining geodynamo in the outer core, while rotation is an essential organizer rather than a standalone source. Its observed motion is a coupled signature of changing core dynamics, external currents, induction in conducting layers, and measurement geometry. The proposed experiment should be accepted only if it improves held-out calibrated prediction and source attribution without hiding non-identifiability. Until then, it is a rigorous design for learning what the field is doing, not a claim that its future has already been computed.

\begin{thebibliography}{00}
\bibitem{robertsking} P. H. Roberts and E. M. King, ``On the genesis of the Earth's magnetism,'' \textit{Reports on Progress in Physics}, vol. 76, no. 9, Art. no. 096801, 2013, doi: 10.1088/0034-4885/76/9/096801.
\bibitem{buffett} B. A. Buffett, ``Earth's core and the geodynamo,'' \textit{Science}, vol. 288, no. 5473, pp. 2007--2012, 2000, doi: 10.1126/science.288.5473.2007.
\bibitem{glatzmaier} G. A. Glatzmaier and P. H. Roberts, ``A three-dimensional self-consistent computer simulation of a geomagnetic field reversal,'' \textit{Physics of the Earth and Planetary Interiors}, vol. 91, nos. 1--3, pp. 63--75, 1995, doi: 10.1016/0031-9201(95)03049-7.
\bibitem{christensen} U. R. Christensen and J. Aubert, ``Scaling properties of convection-driven dynamos in rotating spherical shells and application to planetary magnetic fields,'' \textit{Geophysical Journal International}, vol. 166, no. 1, pp. 97--114, 2006, doi: 10.1111/j.1365-246X.2006.03009.x.
\bibitem{alken} P. Alken \textit{et al.}, ``International Geomagnetic Reference Field: the thirteenth generation,'' \textit{Earth, Planets and Space}, vol. 73, Art. no. 49, 2021, doi: 10.1186/s40623-020-01288-x.
\bibitem{finlay} C. C. Finlay \textit{et al.}, ``Recent geomagnetic secular variation from Swarm and ground observatories as estimated in the CHAOS-6 geomagnetic field model,'' \textit{Earth, Planets and Space}, vol. 68, Art. no. 112, 2016, doi: 10.1186/s40623-016-0486-1.
\end{thebibliography}

\end{document}
